%%==========================================================================%%
%% datarecords.tex -- "Data Records".                                        %%
%% DICOM release (AutoPET-shaped): DICOM hierarchy + DICOM-SEG labels +       %%
%% a per-examination metadata CSV. NO summary statistics here (those are in   %%
%% the Data Properties subsection of Methods). American spelling.            %%
%% Counts, formats, and identifiers are [TBD] pending reconciliation.        %%
%%==========================================================================%%

\section*{Data Records}\label{sec:datarecords}

The dataset is deposited at [TBD: repository] under accession [TBD] and is available at [TBD: Digital Object Identifier], under [TBD: license, pending the open-versus-controlled-access decision].
The imaging, the segmentation labels, and the indexing metadata are organized as described below.
This section describes the files, their formats, the folder structure, and the field definitions.
The lesion-burden distributions are reported in the Data Properties subsection of the Methods, and no summary statistics appear here.

\subsection*{DICOM data}\label{subsec:dicom}
The imaging is released in DICOM, the native clinical format, organized under the standard DICOM hierarchy and indexed by an anonymized subject identifier.
Each subject contains one or more imaging studies (examinations), and each study contains the FDG-PET and MRI series acquired in that session, together with the released segmentation series.
Studies and series are identified by their DICOM unique identifiers, which are remapped during de-identification, and a per-subject date offset is applied so that the relative timing between a subject's examinations is preserved while the absolute dates are removed.
The release comprises [TBD: number] examinations across [TBD: number] subjects, [TBD: number] image series, and [TBD: number] DICOM files, totaling approximately [TBD: size] [TBD: confirm against the DICOM-deliverable manifest, not the full identified cohort].
The directory layout is shown in Figure~\ref{fig:layout}.

%%-- Directory-tree figure placeholder (cf. AutoPET Fig. 2) -----------------%%
\begin{figure}[t]
\centering
\fbox{\parbox[c][5cm][c]{0.85\textwidth}{\centering\itshape
[Figure 2 placeholder: directory tree of the release, showing the DICOM
subject/study/series hierarchy, the DICOM-SEG label series, the derived NIfTI
volumes, and the metadata table. Real image pending, sans-serif, white
background at submission.]}}
\caption{\textbf{Directory layout of the release.}
The DICOM imaging is organized by anonymized subject and study, the tumor-lesion segmentations are stored as DICOM-SEG series aligned to the corresponding PET, and the indexing metadata is provided as a plain-text table.}
\label{fig:layout}
\end{figure}

\subsection*{Segmentation labels}\label{subsec:labels}
The tumor-lesion segmentations are released as DICOM-SEG series, each spatially aligned to the FDG-avid PET series of the same examination, and each recording its target structure and label type (voxel-level segmentation).
Lesion masks are provided for the lesion-positive subset (approximately [TBD] examinations).
The disease-free reference examinations (approximately [TBD] examinations) carry no lesion mask and are flagged in the metadata as complete-metabolic-response follow-up studies, so that they can be selected as pathology-negative references.

\subsection*{Derived formats}\label{subsec:derived}
For convenience, the release includes the conversion utilities that read the DICOM imaging and produce derived NIfTI volumes, apply the SUV conversion to the PET, and resample the PET and MRI onto a common grid where required (Section~\ref{sec:code}).
[TBD: state whether the derived NIfTI volumes are shipped with the release or are regenerated by the user from the DICOM via the conversion utilities.]

\subsection*{Metadata}\label{subsec:metadata}
A per-examination metadata table is provided as an open plain-text file (CSV) with an accompanying data dictionary, so the dataset is usable with standard tooling.
The table captures the indexing variables useful for cohort selection and modeling, grouped as in Table~\ref{tab:vars}, and every non-obvious column is defined in the data dictionary.
The acquisition fields are populated from the DICOM headers.
The release contains no radiology reports and no staging or treatment-response labels.
Lymphoma class and the complete-metabolic-response flag are the only clinical fields provided, and the staging context that motivates the cohort is not part of the released labels.

\begin{table}[h]
\caption{Indexing variables provided with the dataset. Final field names and definitions are [TBD].}\label{tab:vars}
\begin{tabular}{@{}ll@{}}
\toprule
Group & Example fields \\
\midrule
Subject / demographics & subject\_id, age\_at\_scan, sex \\
Clinical               & lymphoma\_class (Hodgkin / non-Hodgkin), subset\_flag (lesion-positive / disease-free reference) \\
Acquisition            & scanner, field\_strength, injected\_activity, uptake\_time, sequence/contrast \\
Labels                 & label\_id, structure, label\_type, segmentation\_reference\_series \\
File / QC              & dicom\_reference, checksum \\
\botrule
\end{tabular}
\end{table}

\subsection*{Related data}\label{subsec:related}
The paired MRI series released here are the same co-acquired series released, with multi-structure anatomical segmentations, in the companion descriptor drawn from this cohort, and approximately [TBD] examinations from [TBD] patients appear in both releases.
The split between the two releases is by labeled task (tumor lesions here, anatomy in the companion descriptor) rather than by imaging modality.
The companion descriptor will be cited as a numbered data citation once it has an assigned Digital Object Identifier [TBD: companion accession and DOI].
The adult whole-body FDG-PET/CT lesion dataset used as the comparator in the Technical Validation is cited as deposited data~\cite{Gatidis_2022_data}.
